\chapter{Quadro storiografico e metodologico}


% ------------------------------------------------------------------
\section{Stato dell'arte}
\label{sec:stato_arte}
% ------------------------------------------------------------------

L'\textit{Ottica delle oscillazioni elettriche} \cite{righi1897ottica} è il trattato in cui Augusto Righi (1850--1920)
racchiude e ordina le esperienze sulle onde elettromagnetiche da lui eseguite per ottenere prove sperimentali che confermassero
la teoria di Maxwell, seguendo la strada già spianata da Heinrich Rudolph Hertz (1857--1894) \cite{hertz1893electric}.\\
Ma quale significato ha il titolo?\\
Il sottotitolo parla di «Studio Sperimentale sulla produzione di fenomeni analoghi ai principali fenomeni ottici»: in che senso, dunque, 
le oscillazioni elettriche possono essere trattate con le leggi dell'ottica? E quale ruolo svolge l'analogia in questo procedimento?\\
Nell'opera monumentale di Dragoni \cite{dragonirighi2023} vengono descritti la biografia, i meriti accademici e il ruolo del fisico italiano
nel panorama internazionale dell'epoca, ma alla questione dell'analogia viene accennata di sfuggita \cite[p.~160]{dragonirighi2023} 
\cite[p.~234--235]{dragonirighi2023}.\\
Bertozzi \cite{bertozzi2020augusto} tratta, sul \textit{Nuovo Saggiatore}, il ruolo dell'intuizione sperimentale negli esperimenti di Righi.
L'analogia compare in alcuni passaggi, ma non diventa un tema di analisi.\\ 
Nel passo tratto da l'Ottica delle oscillazioni elettriche che Bertozzi riporta \cite[p.~55]{bertozzi2020augusto}, Righi osserva che gli esperimenti di Hertz mostrano un'analogia completa 
tra la propagazione dei raggi di forza elettrica e quella dei raggi luminosi, e ne ricava il problema tecnico di
ridurre la lunghezza d'onda; altrove si parla dell'effetto Righi-Leduc come analogo termico dell'effetto Hall
\cite[p.~55]{bertozzi2020augusto} e di doppietti planetari ipotizzati per analogia con le stelle doppie \cite[p.~60]{bertozzi2020augusto}.\\
Per le onde elettromagnetiche, tuttavia, il programma di Righi è descritto come il completamento di quello iniziato da Hertz, volto a mostrare che luce
e oscillazioni elettriche sono entrambe onde elettromagnetiche \cite[pp.~56 e~62]{bertozzi2020augusto}.

\section{L'analogia nella scienza}

Prima di addentrarci in ciò che il titolo di questa sezione suggerisce, è necessario fare una premessa: questa parte non aspira tanto a essere
un trattato filosofico sul concetto di analogia, quanto a dare un
riferimento su che cosa sia l'analogia e quale ruolo svolga nella scienza.\\
Come primo termine di confronto prendiamo la prima definizione che ne dà la Treccani:
«\textit{Rapporto di somiglianza tra due oggetti, tale che dall’eguaglianza o somiglianza constatata
tra alcuni elementi di tali oggetti si possa dedurre l’eguaglianza o somiglianza anche di
tutti gli altri loro elementi}» \cite{analogia_treccani}.\\
Seguendo questa definizione, un esempio di analogia potrebbe essere questo: un flusso di elettroni scorre (anche detto corrente elettrica)
attraverso un conduttore analogamente a come farebbe un liquido in un tubo. In questo caso i due oggetti sono «flusso di elettroni in un conduttore» e
«liquido in un tubo».\\ 
Sfruttando questa somiglianza di comportamento risulta più semplice dedurre come si comporterebbe il flusso di elettroni
qualora dovesse incontrare una strettoia nel conduttore.\\
L'esempio posto poc'anzi è utile per capire il significato di «analogia» riportato dal dizionario, ma mette anche in risalto il ruolo «euristico» \cite{Shelley2019} che
l'analogia acquisisce nella scienza.\\
Per ruolo «euristico» si intende l'uso dell'analogia come aiuto per la ricerca, per favorire la comprensione e per formulare possibili soluzioni
ai problemi \cite{Shelley2019}.\\
Anche Melandri \cite[pp.~22--23]{melandri2023analogia} riconosce all'analogia una funzione euristica, cioè di scoperta di ipotesi, 
ma ne sottolinea la «validità transitoria»: una volta confermata l'ipotesi, il suo valore di verità non dipende più dalla genesi analogica.\\
Tornando alla scienza, riporto un passaggio significativo di Nappo et al. \cite{nappo2026analogical}:

\begin{quote}
    \textit{In che modo gli scienziati portano avanti le loro ricerche? [...], una strategia comune alle
    scienze naturali e alle scienze sociali consiste nel ricorrere all’analogia.}(traduzione personale)
\end{quote}
Apprezzo questo passaggio perché offre una prospettiva più concreta riguardo la ricerca: l'analogia può essere una buona fonte di ispirazione
per trovare nuovi percorsi di ricerca che altrimenti sarebbero stati preclusi.\\

